% ============================================================================
% Chapter 3 — Requirements, Impact, and Project Management
% Rubric-mandated section structure preserved at the top level; depth lives
% in subsections.
%
% Course outcomes covered: CO2 (specifications and standards), CO3
% (societal impact), CO4 (sustainability and environment), CO10
% (resource management), CO11 (economic analysis), CO12 (ethics).
% Section titles 3.1 through 3.8 are fixed by the rubric.
% ============================================================================

\chapter{Requirements, Impact, and Project Management}
\label{ch:requirements}

% Chapter scope: states the functional and non-functional requirements
% the architecture must meet, surveys the societal, environmental, and
% ethical considerations, declares the engineering standards followed,
% and lays out the project management plan, risk register, and
% economic analysis. This chapter operationalises the Course Outcomes
% CO2, CO3, CO4, CO10, CO11, and CO12.

\section{Final Specifications and Requirements}
\label{sec:specifications}

\subsection{Functional requirements}
\label{sec:func-reqs}

The architecture must perform seven functions at deployment, each derived from the problem statement in \Cref{ch:introduction}; \Cref{tab:func-reqs} lists them with the component that delivers each. Three of them need a word of explanation. Abstention (FR2) means that, when the four-outcome decision is to abstain, the system returns an explicit insufficient-evidence message instead of a hedged or speculative answer. Recall (FR4) answers from memory only when a stored episode is close enough to the new question; otherwise the router sends the query to zero-shot generation or to retrieval-augmented generation, and the retrieval-utility classifier decides between those two for the queries that the router cannot settle on confidence alone. Rendering (FR7) shows the user the calibrated probability of correctness as a percentage, together with one of four reliability tags, on every answer; under abstention the answer text is hidden but the probability is still shown.

\begin{table}[!htbp]
\centering
\small
\begin{tabular}{@{}p{0.8cm}p{4.4cm}p{8cm}@{}}
\toprule
\textbf{ID} & \textbf{Requirement} & \textbf{Architectural component (specified in \Cref{ch:methodology})} \\
\midrule
FR1 & Calibrated answer per query & Nine-signal verifier feeding the per-benchmark calibrated probability composite (with shrinkage prior) \\
FR2 & Explicit abstention under uncertainty & Fixed-threshold storage gate $+$ confabulation early-exit rule \\
FR3 & Commit verified episodes to memory as one record (question, answer, supporting context, storage score) & Four-outcome decision tree (store branch) \\
FR4 & Recall on subsequent queries & Three-tier router with combined-score dispatch and safety override; the retrieval-utility classifier re-decides, at threshold $\tau_{\text{RUC}} = 0.375$, the queries that reach the safety override or fall through the score formula \\
FR5 & Cycle-boundary fine-tune on verified episodes only & Self-improvement loop with low-rank adapter, pool-reweighting layer, and strict training threshold \\
FR6 & Retention guard with adapter rollback that keeps the memory & Multi-modal retention probe panel (worst-probe rule) and asymmetric memory-versus-adapter rollback \\
FR7 & Calibrated probability rendered to the user with every answer & Production response envelope: four-class decision tag $+$ continuous $u_{\text{stored}}$ percentage; insufficient-evidence message replaces the answer text under abstain \\
\bottomrule
\end{tabular}
\caption[Functional requirements and the architectural components that deliver them.]{Functional requirements and the architectural components that deliver them.}
\label{tab:func-reqs}
\end{table}

\subsection{Non-functional requirements}
\label{sec:nonfunc-reqs}

Five non-functional requirements bound the operating envelope of the architecture (\Cref{tab:nonfunc-reqs}).

\textbf{Latency:} The cost of a query depends on its tier: the memory-direct tier costs an embedding lookup and a similarity search, the zero-shot tier adds one decode, and the retrieval-augmented tier adds passage retrieval and a grounded decode. Per-tier latencies are reported in \Cref{ch:evaluation}.

\textbf{Storage precision:} Among the answers that the gate stores, the fraction that are correct should be high. The requirement is stated on the calibration fold at the locked cut-point $\tau_{\text{store}} = 0.60$, where the storage precision $\pi_{\text{store}}$ is
\begin{equation}
\Pr\!\big(\mathrm{EM} = 1 \,\big|\, u_{\text{stored}} \ge \tau_{\text{store}}\big)_{\text{cal}} \;=\; \pi_{\text{store}}.
\label{eq:nfr-precision-floor}
\end{equation}
The cut-point was chosen on the cycle-zero calibration fold before the main run, and the realised precision at every cycle is reported in \Cref{ch:evaluation}.

\textbf{Retention:} After every fine-tune, accuracy on each probe of the multi-modal probe panel must stay at or above $93$ percent of its value before any fine-tuning; otherwise the cycle is aborted and the adapter is rolled back:
\begin{equation}
\min_{p \in \mathcal{P}}\; \rho_{p}(M_{t}) \;:=\; \min_{p \in \mathcal{P}}\; \frac{\mathrm{EM}(M_{t},\,p)}{\mathrm{EM}(M_{0},\,p)} \;\ge\; \rho_{\min} = 0.93,
\label{eq:nfr-retention-floor}
\end{equation}
where $M_{t}$ is the generator after the fine-tune at cycle $t$ and $\mathcal{P}$ is the probe panel: a general-knowledge multiple-choice probe (MMLU), an open-text factoid probe drawn from the TriviaQA test fold, and a commonsense multiple-choice probe drawn from the CommonsenseQA test fold.

\textbf{Reproducibility:} Every released artefact, including the locked cut-points, the calibrated composite, the seed-pool composition and the per-cycle snapshots, must be regenerable from one recorded random seed and a published code commit.

\textbf{Integrity:} No sample in the evaluation fold may appear in the calibration fold, the seed pool or the per-cycle stream chunks, so that calibration-fold and evaluation-fold results are measured on disjoint samples.

\begin{table}[!htbp]
\centering
\small
\begin{tabular}{@{}p{0.8cm}p{3.5cm}p{4.5cm}p{4.5cm}@{}}
\toprule
\textbf{ID} & \textbf{Requirement} & \textbf{Target} & \textbf{Where it is checked} \\
\midrule
NFR1 & Per-query latency & Tier-dependent: embedding lookup (Tier~1), single decode (Tier~2), retrieval $+$ grounded decode (Tier~3) & Per-tier latency in \Cref{ch:evaluation} \\
NFR2 & Storage precision at $\tau_{\text{store}} = 0.60$ & Calibration-fold precision of the stored class, with the cut-point fixed before the main run & Per-cycle precision in \Cref{ch:evaluation} \\
NFR3 & Retention floor (multi-modal probe panel) & $\ge 0.93$ of the pre-fine-tune accuracy on every probe; the worst probe decides & Per-cycle probe ratios in \Cref{ch:evaluation}; one breach (cycle four) rolled back \\
NFR4 & Reproducibility & Every artefact regenerable from one seed and a published commit & Locked configuration files and archived snapshots (\Cref{app:code-data-availability}) \\
NFR5 & Integrity & Evaluation, calibration, seed and stream-chunk folds disjoint & Disjointness verified per benchmark \\
\bottomrule
\end{tabular}
\caption[Non-functional requirements, their targets, and where each is checked.]{Non-functional requirements, their targets, and where each is checked.}
\label{tab:nonfunc-reqs}
\end{table}

\subsection{Constraints and design commitments}
\label{sec:constraints}

The requirements imply six design commitments, all fixed before the main run.
\begin{itemize}
    \item \textbf{Single backbone:} The generator backbone is the same in every cycle and every tier; only the low-rank adapter is trained, so reported gains come from memory, calibration and adapter fine-tuning rather than from a change of backbone.
    \item \textbf{Frozen cut-points:} The storage and deferral cut-points, $\tau_{\text{store}} = 0.60$ and $\tau_{\text{defer}} = 0.45$ on the calibrated probability scale, are fixed once before the main run and never changed, so the operating point of the storage decision does not drift. Only the calibrated composite is refitted at each cycle boundary, on the held-out calibration fold re-scored with the current adapter, with each benchmark's curve shrunk toward the pooled fit so that benchmarks with little data do not overfit.
    \item \textbf{Local execution:} No external model service sits in the inference path: the generator, the verifier ensemble, the embedding encoder and the passage index all run from open weights on one rented graphics processor. The language-model judge that scores TruthfulQA belongs to the evaluation, not to the system.
    \item \textbf{Bounded cycles:} The self-improvement loop is capped at ten cycles and may stop earlier on an automatic early-stop check (\Cref{ch:methodology}).
    \item \textbf{Matched protocol:} Every external baseline is evaluated on the same samples, prompts and scoring rules as the architecture, so that the comparison in \Cref{ch:evaluation} is not confounded by re-implementation choices.
    \item \textbf{Single seed:} All randomness derives from one recorded seed, and the per-cycle artefacts are archived so that a third party can regenerate the trajectory.
\end{itemize}

\section{Societal Impact}
\label{sec:societal-impact}

\subsection{Deployment benefits}
\label{sec:soc-benefits}

The main societal benefit of the architecture is a reduction in confident misinformation from deployed language models on factual question answering. Every answer carries a calibrated probability and a reliability tag, and the answer is withheld when the evidence is insufficient. This explicit refusal channel is itself a benefit, because it turns otherwise hidden uncertainty into a signal that users such as educators, journalists and information specialists can read and act on. Fluent but ungrounded answers have caused documented harm in health information: Alkaissi and McFarlane~\cite{alkaissi2023artificial} report fabricated clinical references, and Chelli et al.~\cite{info:doi/10.2196/53164} report high rates of hallucinated references in systematic-review searches. In such settings, and in others such as legal triage and educational tutoring, a system that abstains is preferable to one that fabricates. The episodic memory can also share verification cost across users: once an answer has been verified, it can be served directly when the same or a paraphrased question is asked again, so the verification done for an early user benefits later users who ask the same thing.

\subsection{Dual-use considerations}
\label{sec:soc-dual-use}

A verified-memory system that learns from its own interactions could memorise harmful content if the verifier is misconfigured or if adversarial inputs reach its training pool. Three mechanisms limit this risk.

\textbf{Retroactive re-verification:} At each cycle boundary every stored episode is re-scored with the updated generator and the refitted composite; its score can go down as well as up, and entries whose score falls below $0.50$ are removed, so a bad entry does not persist unchallenged.

\textbf{Strict training threshold:} The training pool admits only episodes whose storage score is at least $0.70$, above the storage cut-point of $0.60$, so marginal episodes stay out of the fine-tune even when they are admissible for retrieval.

\textbf{Retention guard:} If a fine-tune degrades any probe in the multi-modal panel (general-knowledge multiple choice plus held-out test slices of two training benchmarks) below the floor, the cycle is aborted and the adapter is rolled back, so one damaging cycle cannot silently carry forward.

These mechanisms reduce the risk but do not remove it; how far they reach depends on the verifier's coverage, which is reported empirically and discussed as a threat to validity in \Cref{ch:evaluation}.

\subsection{Marginalised-user impact}
\label{sec:soc-marginalised}

The harm caused by confident hallucination is not evenly distributed. Users who cannot independently check a confident-sounding answer, such as students using language-model tutors, patients using language-model health agents, and applicants using language-model administrative agents, are the most exposed, because they have no domain expert to fall back on. The calibrated abstention class is a partial response: by returning an explicit refusal rather than a hedged guess when confidence is low and little supporting evidence is found, the system reduces the most dangerous failure mode for exactly these users. The abstention is governed by a fixed cut-point on a calibrated probability rather than by a soft heuristic, and the abstention rate is reported in the four-outcome decision distribution in \Cref{ch:evaluation}. The architecture does not solve unequal access to verification, but it narrows one channel through which fluency is mistaken for reliability.

\section{Environmental Impact}
\label{sec:environment}

\subsection{Compute footprint of the experiment}
\label{sec:env-experiment}

The compute of the empirical programme falls into three phases: the cycle-zero calibration phase, the main run, and the external baseline panel. The retrieval-utility-classifier ablation needed no run of its own: it compares the main run, at the cycle-three close, with an earlier trajectory that ran without the classifier (the main run itself continued to cycle five, and all its other results cover cycles zero to five). Everything ran on a single rented graphics processor of the consumer $32$-gigabyte class (an RTX 5090). The main run was the largest: according to its run log, cycles zero to five took $6.0$, $29.0$, $35.4$, $43.2$, $24.3$ and $54.6$ hours, about $193$ hours (eight days) in total. Part of the growth comes from the retroactive re-verification pass, which grows with the memory (about $38$ minutes for $431$ entries at cycle one, about $16$ hours for $10{,}214$ entries at cycle five). The main run was not the first. Three earlier trajectories were rejected or superseded: the first benchmark panel, a run stopped in May 2026, and the trajectory without the classifier. The last took about $134$ hours over cycles zero to five according to its run log, and its first three cycles supply the classifier-off arm of the ablation; the run logs of the first two were not archived, so their compute cannot be reported exactly. Counting these runs, the programme used several times the compute of the main run alone. The carbon footprint of a phase can be estimated with the method of \textcite{strubell2019energy},
\begin{equation}
\mathrm{kgCO_{2}e}_{\,\phi} \;=\; T_{\phi} \cdot P_{\mathrm{TDP}} \cdot \varepsilon_{\mathrm{grid}},
\label{eq:carbon-per-phase}
\end{equation}
where $T_{\phi}$ is the wall-clock duration of phase $\phi$ in hours, $P_{\mathrm{TDP}}$ is the rated power of the graphics processor in kilowatts, and $\varepsilon_{\mathrm{grid}}$ is the local grid emissions factor in kilograms of carbon-dioxide equivalent per kilowatt-hour. At the card's rated board power of $0.575$ kilowatts, the main run used at most about $111$ kilowatt-hours of graphics-processor energy. The grid factor of the rented host is not known, so no carbon figure is reported.

\subsection{Cycle budget rationale}
\label{sec:env-cycles}

The ten-cycle cap bounds the compute that the loop may spend. The convergence heuristic of \Cref{ch:methodology} (\Cref{thm:convergence}; a heuristic, not a proven result) suggests that gains flatten within a few cycles, and Song et al.~\cite{song2024gap} find that iterative self-improvement saturates within a few rounds; a cap of ten leaves room above that while bounding cost. Each cycle adds its own fine-tuning and re-verification energy, so stopping once gains have flattened is an environmental choice as well as an experimental one. The realised run stopped after cycle five, as planned once cycle four had been aborted by the retention guard and its weights rolled back. The automatic early-stop check fired at the same point, but two of its three inputs were wired to the wrong quantities: its storage-saturation signal read the storage rate of the evaluation fold, where nothing is ever stored, instead of the stream admission rate (about half of each chunk), and its retention signal compared raw MMLU accuracy (about $0.61$) with the $0.93$ ratio floor instead of the retention ratio (between $0.94$ and $1.00$). Both signals were therefore certain to fire as soon as the check first ran, at cycle five, so its firing is not evidence that the trajectory had levelled off. The wiring was corrected after the run; replaying the check on the realised cycle values with the corrected inputs, none of its three signals fires at cycle five. The trajectory is reported in \Cref{ch:evaluation}.

\subsection{Retention versus retraining}
\label{sec:env-retention}

The retention guard with rollback also saves compute. Without it, a loop that regressed would either keep training from damaged weights, making every later cycle pay to recover the lost capability, or need a restart from a clean checkpoint, discarding all the fine-tuning done since. With it, a failed cycle costs only its own fine-tuning: the adapter returns to the previous checkpoint, while the memory, including the episodes verified during the failed cycle, is kept. The argument is qualitative; the realised run had one such rollback, at cycle four.

\section{Ethical Issues}
\label{sec:ethics}

\subsection{Alignment and refusal posture}
\label{sec:ethics-alignment}

The architecture treats explicit abstention as a first-class action rather than a fallback, and this choice has an ethical justification beyond its effect on accuracy. A model that hedges or speculates under uncertainty exposes the user to harm without telling them that uncertainty is present, whereas a model that abstains hands the uncertainty to the user as a visible signal, on which they can act by seeking a second opinion or a human expert. This connects to the question, studied by Kadavath et al.~\cite{kadavath2022language}, of whether language models know what they know. The storage and deferral cut-points are fixed values on the calibrated probability scale, chosen on the cycle-zero calibration fold before the main run rather than tuned by hand during it. The ethical commitment behind this posture is that a model which cannot be confident should not present its answer as if it were.

\subsection{Bias propagation through episodic memory}
\label{sec:ethics-bias}

A system that commits its own outputs to memory and then trains on that memory can amplify biases in the verifier or in the underlying data. Three mechanisms limit this. The retroactive re-verification pass re-scores every stored episode at every cycle boundary and removes entries whose score falls below $0.50$, so a bias that enters the stored pool does not persist unchallenged. The per-cycle refit of the calibrated composite, with each benchmark's curve shrunk toward the pooled fit, follows score-distribution drift by re-scoring the held-out calibration fold with each cycle's adapter, while the storage cut-point stays fixed so that the gate's operating point does not move. The strict training threshold holds the training pool to a higher bar than the storage tier (a storage score of at least $0.70$ against $0.60$ for storage), so marginal episodes cannot enter the next fine-tune even when they pass the retrieval gate. None of these mechanisms removes bias inherited from the pre-training corpus, which remains a limitation of the approach; what they do is limit the rate at which new bias can enter through the self-improvement loop itself.

\subsection{Data provenance and licensing}
\label{sec:ethics-provenance}

The empirical programme uses publicly released benchmarks and a public retrieval corpus throughout. The five evaluation benchmarks are FEVER \cite{thorne-etal-2018-fever}, TriviaQA \cite{joshi-etal-2017-triviaqa}, CommonsenseQA \cite{talmor-etal-2019-commonsenseqa}, TruthfulQA \cite{lin-etal-2022-truthfulqa}, and StrategyQA \cite{10.1162/tacl_a_00370}, all distributed under licences that permit research use with attribution. Two further benchmarks, HotpotQA and Natural Questions, were evaluated at cycle zero and dropped before the main run (\Cref{sec:scope-research}); their cycle-zero results are kept in the project archive. The retrieval corpus is a Wikipedia-derived passage collection distributed under a Creative Commons licence. No private user data, scraped social-media content or proprietary corpus enters the calibration fold, the seed pool or the cycle pool. The cold-start memory that seeds the self-improvement loop is built from the benchmarks' training splits, with disjointness from the evaluation fold enforced at the sample level. The code is public, and the configuration files, the per-cycle artefacts and the pre-self-improvement snapshot are archived in a shared folder, so that the data path can be audited end to end (\Cref{app:code-data-availability}).

\section{Standards and Codes of Practice}
\label{sec:standards}

\subsection{Software-engineering standards}
\label{sec:std-software}

The implementation is written in Python and broadly follows the PEP~8 conventions for layout and naming. The project does not ship a pinned requirements file; the versions of the main libraries in the realised environment are recorded in \Cref{app:repro-versions}. All randomness is seeded from one value at process start (the Python, NumPy and PyTorch generators), and deterministic cuDNN kernels are selected where available; the attention kernels on the graphics processor are not bitwise deterministic, so repeated runs agree up to floating-point noise rather than bit for bit. The source is kept under Git version control in a public repository, with one branch per project phase, and each cycle's artefacts are archived at the cycle boundary, so that the empirical record can be traced to the code that produced it.

\subsection{Reproducibility standards}
\label{sec:std-reproducibility}

Reproducibility is treated as a deliverable rather than a courtesy. Four kinds of artefact support it.
\begin{itemize}
    \item \textbf{Recipe:} \Cref{app:reproducibility} gives the run sequence from data preparation to the final tables, together with a checklist for comparing reproduced numbers against the thesis.
    \item \textbf{Pre-self-improvement snapshot:} The cold-start memory, the cycle-zero calibrated composite and the configuration at the start of the main run are archived, so that a run can resume from that point instead of from scratch.
    \item \textbf{Per-cycle artefacts:} Each cycle's run summary, memory snapshot, retention-probe readings, precision audit and decision log are kept together.
    \item \textbf{Locked configuration:} One configuration file records the storage and deferral cut-points, the shrinkage prior, the boost-layer regulariser and the low-rank adapter settings used at every downstream scoring and training step.
\end{itemize}
Together these let a reader check a specific claim in \Cref{ch:evaluation} by re-running the stage that produced it rather than by trusting the reported number.

\subsection{Evaluation-protocol standards}
\label{sec:std-evaluation}

The empirical comparisons follow four protocol standards.

\textbf{Matched protocol:} Every external baseline is evaluated on the same samples, the same prompt template family and the same scoring rules as the architecture.

\textbf{Pre-specification:} The calibration sweep grid, the cut-point targets and the expected effect bands of the ablation were recorded in the runbook before the corresponding results were read, and the locked configuration was applied unchanged in every cycle of the main run; outcomes outside the expected bands are reported rather than tuned away. Throughout the thesis, a quantity described as pre-specified or registered was fixed in the project's configuration and runbook before the main evaluation run; the term does not refer to an entry in an external pre-registration registry.

\textbf{Selective prediction:} Abstentions are reported per benchmark alongside accuracy, following the selective-prediction framing of Kamath et al.~\cite{kamath2020selective}, in which a system that may abstain is judged both on how often it answers and on how accurate its answers are.

\textbf{Multiple comparisons:} Within each benchmark, the baseline contrasts are corrected with the Holm--Bonferroni procedure \cite{holm1979simple}, so that a significance verdict cannot be inflated by the number of comparisons. Each contrast uses the paired McNemar test with the continuity correction of Edwards~\cite{edwards1948note}, and a bootstrap $95$ percent confidence interval on the exact-match difference is reported alongside the corrected $p$-value, so that effect size is visible as well as significance.

\section{Project Management Plan}
\label{sec:project-management}

\subsection{Phase breakdown}
\label{sec:pm-phases}

The project has five phases, each producing what the next consumes.
\begin{enumerate}
    \item \textbf{Pre-thesis:} Literature review, problem statement, methodology design, and the per-stage pipeline that the empirical phases depend on.
    \item \textbf{Calibration:} Fit the calibrated composite on the disjoint cycle-zero calibration fold, fix the storage and deferral cut-points, and lock the configuration for the main run (only the composite is refitted in later cycles).
    \item \textbf{Main run:} Iterate the per-query pipeline and the cycle-boundary passes for up to ten cycles, producing the per-cycle artefacts.
    \item \textbf{Comparison:} Run the external baseline panel and the retrieval-utility-classifier on/off ablation under matched protocol, followed by the statistical analysis.
    \item \textbf{Write-up:} The thesis report, the defence presentation, and a journal manuscript planned for after the defence.
\end{enumerate}

\subsection{Milestone schedule}
\label{sec:pm-milestones}

The pre-thesis milestone, a locked methodology design, was met on schedule. The calibration milestone, the cycle-zero validation gate, took two iterations: the first failed because of a bug in the step that writes out the verifier signals, and the second, with the bug fixed and a wider parameter grid, produced the locked configuration that the main run uses. The slip is reported here rather than smoothed over, and \Cref{ch:evaluation} gives the diagnostics of both iterations. The main run completed six cycles (cycles zero to five) on 27 May 2026, and the comparison phase completed the eight-baseline panel and the retrieval-utility-classifier ablation. The write-up milestone is the submission of this report and its companion artefacts.

\subsection{Deliverables}
\label{sec:pm-deliverables}

The project releases five deliverables.
\begin{itemize}
    \item \textbf{Source code:} public at \url{https://github.com/aksaN000/caem-thesis}: the implementation of the architecture, the cycle-zero calibration scripts, the cycle-boundary passes, and the baseline and ablation drivers.
    \item \textbf{Artefact bundle:} stored in the shared Google Drive folder \url{https://drive.google.com/drive/folders/1xErhsRDYe_qO0ZgUBn0WXK81eHc5jZCa}: the cycle-zero calibrated composite, the locked cut-points, the cold-start memory snapshot, and the per-cycle artefacts of the realised trajectory. The $64$-gigabyte Wikipedia passage index is mirrored on a private Hugging Face repository, available on request (\Cref{app:code-data-availability}).
    \item \textbf{Thesis report:} this document.
    \item \textbf{Defence presentation:} a talk summarising the headline result, the method and the threats to validity.
    \item \textbf{Journal manuscript:} planned for after the defence and prepared from the evaluation chapter.
\end{itemize}

\section{Risk Management}
\label{sec:risk}

\subsection{Technical risks}
\label{sec:risk-technical}

Six technical risks were identified before the main run. Each is paired with its mitigation below, and \Cref{tab:risk-register} records whether it materialised.
\begin{itemize}
    \item \textbf{Calibration drift:} The composite, refitted every cycle, could lose its precision at the frozen cut-point as each adapter update shifts the score distribution. Mitigation: the per-cycle refit (per-signal isotonic curves, per-benchmark shrinkage and boost layer) and a per-cycle precision audit at the locked cut-points.
    \item \textbf{Verifier and judge disagreement:} The short-answer entailment judge and the long-answer judge could score the same answer inconsistently. Mitigation: drop the long-answer judge if its agreement with the short-answer judge misses a fixed correlation floor.
    \item \textbf{Out-of-memory failures:} A long answer could send the verifier down a branch that exceeds the memory of the graphics processor and halt the cycle. Mitigation: chunking the long-answer batch, and a clean restart from the previous checkpoint after an out-of-memory abort.
    \item \textbf{Forgetting:} A fine-tune could erode general capability. Mitigation: the retention guard with adapter rollback (\Cref{eq:nfr-retention-floor}).
    \item \textbf{Early-stop misfire:} The automatic early-stop check could stop the run too early or fail to stop a run that is not converging. Mitigation: requiring two of its three signals to agree, and the ten-cycle cap. In the realised run two of the signals were wired to the wrong quantities, which made the check fire at its first run; the fault is described in \Cref{sec:env-cycles} and was corrected after the run.
    \item \textbf{Retrieval-utility classifier shift:} The classifier is trained on data that deliberately includes a benchmark outside the evaluation panel (the Natural Questions training split), so it may send some queries to the wrong tier on the panel. Mitigation: nested cross-validation, a threshold sweep, and the on/off ablation on the same trajectory.
\end{itemize}

\subsection{Budget and infrastructure risks}
\label{sec:risk-budget}

Three infrastructure risks were identified.
\begin{itemize}
    \item \textbf{Credit exhaustion:} Running out of rented graphics-processor credit before the trajectory and the comparisons finish would leave a partial run that cannot be compared against the baselines. Mitigation: budgeting each phase against the recorded credit, and archiving every state before it is overwritten.
    \item \textbf{Node reclamation:} The provider could reclaim the rented machine in the middle of a cycle. Mitigation: committing artefacts at every cycle boundary, so that a lost machine costs at most the cycle in flight, and the pre-self-improvement snapshot, from which the run can resume.
    \item \textbf{Storage exhaustion:} The artefacts, model snapshots and passage index could exceed the disk of the rented machine or the quota of the hosts that hold them. Mitigation: sizing the artefact catalogue in advance and keeping the larger snapshots and the passage index on separate storage.
\end{itemize}

\subsection{Mitigation register}
\label{sec:risk-mitigation}

\Cref{tab:risk-register} consolidates the register. Its last column records what happened over the completed programme.

\begin{table}[!htbp]
\centering
\footnotesize
\begin{tabular}{@{}>{\raggedright\arraybackslash}p{2.9cm}>{\raggedright\arraybackslash}p{3.9cm}>{\raggedright\arraybackslash}p{4.1cm}>{\raggedright\arraybackslash}p{3.6cm}@{}}
\toprule
\textbf{Risk} & \textbf{Mechanism} & \textbf{Mitigation} & \textbf{Realised} \\
\midrule
Calibration drift across cycles & Composite loses its precision at the frozen $\tau_{\text{store}}$ under adapter-induced shift & Per-cycle composite refit (per-signal isotonic, per-benchmark shrinkage, boost); per-cycle precision audit & \textbf{Partly}: calibration-fold precision held between $0.732$ and $0.746$ over cycles one to five; the calibration-to-evaluation gap is reported in \Cref{ch:evaluation} \\
Verifier and judge disagreement & Short- and long-answer judges score the same answer inconsistently & Drop the long-answer judge if it misses the correlation floor & \textbf{Yes}: long-answer judge dropped (Pearson $0.584 < 0.70$) \\
Out-of-memory failure on long paths & A long answer sends the verifier past the GPU's memory & Chunked long-answer batch; clean restart after an abort & \textbf{Yes}: chunking added; degenerate samples skipped on out-of-memory \\
Forgetting & A fine-tune erodes general capability & Retention guard with adapter rollback & \textbf{Yes}: cycle four, TriviaQA-test ratio $0.886 < 0.93$; adapter rolled back, memory kept \\
Early-stop misfire & Premature stop, or a run that never stops & Two-of-three signal rule; ten-cycle cap & \textbf{Yes}: two inputs mis-wired, so the check fired at its first run (cycle five) regardless of the trajectory; the stop coincided with the planned end; wiring corrected after the run (\Cref{sec:env-cycles}) \\
Classifier distribution shift & Classifier routes some panel queries to the wrong tier & Nested cross-validation; threshold sweep; on/off ablation & Reported in the ablation (\Cref{ch:evaluation}) \\
GPU credit exhaustion & Run does not finish within the rented credit & Phase budgeting; archive before overwrite & \textbf{Yes}, once (May 2026, during an earlier run); resumed from the recorded step after a top-up \\
Node reclamation & Cycle in flight lost when the machine is reclaimed & Per-cycle artefact commit; pre-self-improvement snapshot & Not observed \\
Storage exhaustion & Artefacts and snapshots exceed the available disk & Advance sizing; separate storage for large snapshots & \textbf{Yes}: the rented node's disk filled; the data were copied to a new node with the provider's support and the run resumed from its last checkpoint \\
Calibration iteration failure & Cycle-zero composite fails the validation gate & Re-iterate with the signal-dump bug fixed and a wider grid & \textbf{Yes}: first iteration failed on the signal-dump bug; second iteration locked \\
\bottomrule
\end{tabular}
\caption[Risk register with mitigations and the realised outcome.]{Risk register with mitigations and the realised outcome over the completed programme.}
\label{tab:risk-register}
\end{table}

\section{Economic Analysis}
\label{sec:economic}

\subsection{Compute budget}
\label{sec:econ-budget}

The compute budget is measured in graphics-processor hours on the rented $32$-gigabyte machine used throughout. The main run took about $193$ hours (\Cref{sec:env-experiment}). The cost of the external baseline panel scales with the number of samples and the number of decode passes per sample, with the chain-of-thought and retrieval-augmented baselines the most expensive. The retrieval-utility-classifier ablation added no run: it compares the main run with the earlier trajectory without the classifier at the cycle-three close, the last cycle that both trajectories committed (the earlier one aborted its cycles four and five under the retention guard). \Cref{app:repro-hardware} also gives a reduced configuration that reproduces the trajectory on a $16$-gigabyte device, so that the cost on smaller hardware can be estimated without re-running the experiment.

\subsection{Cost-benefit on the storage gate}
\label{sec:econ-cost-benefit}

On the main run the storage gate admitted about half of each $4{,}700$-sample stream chunk: between $2{,}343$ and $2{,}484$ entries per cycle over cycles one to five (\Cref{tab:app-per-cycle-memory}). Storing an answer adds no verification work at query time, because every signal the gate uses has already been computed for that answer. It does add work later: each stored entry is re-verified at every later cycle boundary, which is why the retroactive pass grows with the memory (\Cref{sec:env-experiment}). The benefit of a stored entry is the decode it saves each time it answers a later question directly,
\begin{equation}
\Delta B_{\text{store}} \;\ge\; n_{\text{hit}} \cdot c_{\text{gen}},
\label{eq:storage-cost-benefit}
\end{equation}
where $n_{\text{hit}}$ is the number of later queries that the entry answers through the memory-direct tier and $c_{\text{gen}}$ is the cost of the zero-shot or retrieval-augmented decode it replaces. This benefit appears only when questions recur, and in the realised run they never did: each cycle drew a fresh stream chunk that shares no question with earlier chunks, and the evaluation fold is disjoint from both, so almost no query matched a stored episode closely enough for the memory-direct tier to answer it. On the evaluation fold that tier answered only $7$ of $9{,}000$ queries. A direct probe of the router (\Cref{ch:evaluation}) shows the tier working when questions do recur: it answered $42$ of $50$ exact repeats and $19$ of $50$ paraphrases of stored questions. The precision of the gate is what makes the reused answers worth serving.

\subsection{Deployment cost projection}
\label{sec:econ-deployment}

By design, the cost per query should fall as the memory grows, because a memory-direct answer costs an embedding lookup and a similarity search instead of a decode. The cost per query at cycle $t$ is a weighted average over the three tiers,
\begin{equation}
\bar{c}(t) \;=\; \pi_{1}(t)\, c_{1} \;+\; \pi_{2}(t)\, c_{2} \;+\; \pi_{3}(t)\, c_{3},
\qquad \pi_{1}(t) + \pi_{2}(t) + \pi_{3}(t) = 1,
\label{eq:deployment-cost}
\end{equation}
where $\pi_{i}(t)$ is the share of incoming queries dispatched to tier $i$ at cycle $t$ and $c_{i}$ is the average per-query cost on that tier. The design expects $\pi_{1}(t)$ to rise and $\pi_{3}(t)$ to fall as cycles progress. In the realised run neither moved as designed: the memory-direct share stayed at or below $0.1$ percent, because by design no stream or evaluation question repeated a stored one (\Cref{sec:econ-cost-benefit}), and the zero-shot and retrieval-augmented shares each stayed between about $45$ and $55$ percent across cycles (\Cref{ch:evaluation}). The memory-direct saving is therefore expected only where real traffic repeats or paraphrases earlier questions, which is the setting that the deployment-study future-work item in \Cref{ch:conclusion} would test under live user traffic.

\section{Chapter signpost}
\label{sec:ch3-signpost}

The requirements, design commitments, societal, environmental and ethical readings, standards, and risk and economic analyses in this chapter fix the constraints under which the architecture is specified. \Cref{ch:methodology} gives that specification in full: the eight-stage per-query pipeline with the Stage 3b retrieval-utility classifier, the nine-signal verifier ensemble and its calibrated probability composite, the fixed-threshold storage gate and the four-outcome decision tree, the cycle-boundary update passes, the constrained self-improvement loop under the multi-modal retention guard, and the pool-purity theorem, the purification condition, and the trajectory heuristics with their proofs or argument sketches and pre-specified empirical receipts.
